\documentclass{article}
\usepackage{amsmath,amssymb,bm,siunitx}
\newcommand{\R}{\mathbb{R}}
\newcommand{\norm}[1]{\left\lVert #1 \right\rVert}
\newcommand{\pd}[3][]{\frac{\partial^{#1} #2}{\partial #3^{#1}}}
\DeclareMathOperator{\Tr}{Tr}
\DeclareMathOperator*{\argmax}{arg\,max}
\def\vec#1{\bm{#1}}

\begin{document}
\section{예제}
벡터 $\vec{x} \in \R^n$ 의 노름은 \(\norm{\vec{x}}\) 이다. % 주석은 제거됨
가격은 \$5 이고, 속도는 $\SI{3.0}{\meter\per\second}$ 이다.

\begin{equation}
  \Tr(A) = \sum_{i} a_{ii} \label{eq:trace}
\end{equation}

식 \eqref{eq:trace} 에서
\begin{align}
  f(x) &= \pd{u}{x} + \pd[2]{u}{t} \label{eq:a} \\
  \hat\theta &= \argmax_{\theta} \mbox{ } L(\theta) \nonumber
\end{align}

\begin{eqnarray*}
  a &=& b + c \\
    &\le& \mathbbm{1} \cdot \sqrt{2}
\end{eqnarray*}

\begin{multline}
  x = a + b + c \\
  + d + e \tag{M}
\end{multline}

\[ \begin{split} y &= \ensuremath{\upalpha} x \\ &= \sfrac{1}{2} \end{split} \]

\verb|$not math$| 그리고 \begin{verbatim}
$ also not math $
\end{verbatim}
\end{document}
